\section{本章小结}
到这里，我们从一条字节流走到了完整消息：先分清角色和设备所有者，再按长度或帧头校验组帧，桥接时统一单位。状态使用有界最新值槽，事件使用确认、重试与去重；连接代次把每轮的数据和清理对应起来。应熟练复现半帧、粘包、坏 CRC、迟到清理和过期数据；真实协议、物理串口与跨机访问范围，按对应设备和当年文档继续核对。

\Needspace{12\baselineskip}
\begin{problemset}
\renewcommand{\theenumi}{\arabic{enumi}}
\item \label{prob:v2-comm-1}基础：同一设备字节流能否由两个互不协调的进程同时解析？建议如何组织读取和分发？
\item \label{prob:v2-comm-2}基础：服务端 bind、listen 后，客户端如何建立连接？随后服务端通过哪个对象收发？
\item \label{prob:v2-comm-3}提高：长度前缀为 3，却只收到两个载荷字节时连接关闭。能否把这半帧留给新连接补完？为什么？
\item \label{prob:v2-comm-4}基础：CRC 失败后为什么不能直接交付载荷？CRC 是否验证发送者身份？
\item \label{prob:v2-comm-5}基础：外部速度 3000 mm/s，内部接口使用 m/s，填多少？还需核对哪些时间与坐标约定？
\item \label{prob:v2-comm-6}提高：旧连接的清理回调在新连接建立后才到达。为什么“先检查代次，再不加锁地关闭当前 socket”仍可能有问题？
\item \label{prob:v2-comm-7}基础：同一位置槽写入 100、200、300 后消费一次，得到什么？若改为三次必须执行的购买事件，组织方法有何不同？
\item \label{prob:v2-comm-8}实践：运行 event\_demo，再模拟第一次确认丢失与所有确认丢失。分别记录实际执行数、发送次数和发送端结果；能否仅凭发送端失败断言业务没有执行？
\item \label{prob:v2-comm-9}提高：当前单调时钟为 12 秒，接收时间为 9 秒，最大年龄为 2 秒，应如何处理？能否直接拿远端未经同步的墙上时钟与本机单调时钟相减？
\end{problemset}

\section{本章验收}
\begin{enumerate}
\item 在两个终端运行 TCP 回显，按检查者指定端口断开并再次连接，展示新代次仍回显、旧清理被拒绝；说明每轮解析器归属。
\item 向解析器注入分片、噪声、坏 CRC 和有效恢复帧，展示只交付完整有效载荷；用 pySerial 伪终端复现读取和关闭异常，说明尚需实机检查的部分。
\item 暂停一条模拟链路，保持另一条继续消费；记录槽数、代次和过期计数。由检查者改过期阈值，解释结果变化。
\item 由检查者更改桥接距离，展示 ROS 话题中的米值；再演示确认丢失与有限重试，区分状态最新值和必须执行的事件。
\end{enumerate}
